% ===================================================================
%  TAVOLA N. 16 — Il grande magazzino. --- Il bazar. --- I giocattoli.
%  Traduction 2026 d'apres texto/io, controlee sur texto/fr et
%  texto/en. Consigne : voir texto/it/10-tavola-01.tex.
%
%  La carte du ciel (bloc 15-1) porte des renvois a LETTRES, (a) a
%  (f), qui se rangent sous le numero 85. L'ido en compose un avec une
%  seule parenthese, « d) » ; l'italien met les deux.
%
%  LA ROSE DES VENTS DE L'ALINEA 15 EST UNE CARTE DE LA MEDITERRANEE
%  DESSINEE DEPUIS UN SEUL POINT.
%
%  Trois colonnes ont deja regarde cette rose. L'allemande a trouve
%  Norden, Süden, Osten, Westen — les mots des marins germaniques, que
%  toute l'Europe a pris. L'estonienne les a REFUSES et nomme ses
%  quatre points par le soleil et par l'heure du repas : põhi le fond,
%  lõuna midi, ida le lever, lääs le coucher. La romanche avait pris
%  ceux du voisin.
%
%  L'italien a les deux jeux. Celui de l'ecole est germanique comme
%  partout — nord, sud, est,
%  ovest —, et c'est celui que la colonne ecrit, parce que
%  c'est celui de la planche. Mais le marin italien en emploie un
%  autre, et il est unique :
%     tramontana  le vent d'au-dela des monts, du nord ;
%     grecale     le vent qui vient de la Grece ;
%     levante     le vent du levant ;
%     scirocco    l'arabe « sharq », l'est — la Syrie ;
%     ostro       le vent du sud, « auster » ;
%     libeccio    l'arabe « lībī », le vent de Libye ;
%     ponente     le vent du couchant ;
%     maestrale   le vent maitre, celui de Rome.
%
%  QUATRE DE CES HUIT NOMS SONT DES NOMS DE PAYS, et si l'on cherche
%  le point d'ou la Grece est au nord-est, la Libye au sud-ouest, la
%  Syrie au sud-est et Rome au nord-ouest, on tombe sur le detroit de
%  Sicile, quelque part entre Malte et le cap Bon.
%
%     CETTE ROSE DES VENTS N'EST PAS UNE ABSTRACTION : C'EST UNE
%     POSITION. Elle a ete faite par des gens qui etaient tous au meme
%     endroit, et elle donne leurs coordonnees.
%
%  La rose germanique dit ou est le soleil ; la rose estonienne dit
%  quelle heure il est ; la rose italienne dit OU L'ON SE TROUVE. Trois
%  facons de s'orienter, et la troisieme est la seule qui ne marche
%  qu'a un seul endroit du monde.
%
%  ET LA MENAGERIE DE CARTON DE L'ALINEA 3 CONFIRME CE QUE LE TABLEAU
%  16 ALLEMAND AVAIT ETABLI. L'allemand composait — Nashorn,
%  Flusspferd, Schildkröte, Fledermaus — et l'estonien avait calque
%  ces composes morpheme pour morpheme. L'italien ne compose rien : il
%  garde les mots grecs et latins des naturalistes,
%  rinoceronte, ippopotamo, coccodrillo,
%  dromedario, pantera, iena. Ce sont les
%  mots de Pline, et ils n'ont pas bouge depuis. L'etiquette de
%  Nuremberg n'avait rien a lui apprendre.
%
%  ET C'EST LA FIN DE LA COLONNE ITALIENNE. Elle aura montre que l'ido
%  NOMME comme l'italien, la ou l'allemande avait montre qu'il COMPOSE
%  comme l'allemand ; que l'Academie ido a delibere deux fois sur des
%  mots qui etaient dans sa propre liste (tableaux 5 et 8) ; et
%  qu'une langue de trois mille ans ne forge pas et n'emprunte guere,
%  parce qu'elle DEPLACE — calcio, giostra, arena, stazione, quatre
%  fois, et c'est sa maniere.
% ===================================================================

%%K t16-apar-1 apar
\VUsaut{4.0mm}
\VUtitre{15.0pt}{60}{TAVOLA N. 16}
\VUsaut{2.0mm}
\VUpk{13.2pt}{40}{Il grande magazzino. --- Il bazar.}
\VUpk{13.2pt}{40}{I giocattoli.}
\VUsaut{3.4mm}

%%K t16-01-1 p
1. --- Il nuovo anno si avvicina. I grandi negozi hanno preparato le loro esposizioni
di \VUgras{giocattoli} \textsuperscript{(1)} e di regali di Capodanno.

%%K t16-01-2 p
Ho chiesto a mio nonno di portarmi al bazar a vedere questa bella esposizione, e lui
ha accettato.

%%K t16-02-1 p
2. --- Prima ci siamo fermati davanti a certi bei trofei d'armi con un \VUgras{fucile},
un \VUgras{berretto d'ordinanza}, una \VUgras{sciabola}, un \VUgras{cinturone} e un
paio di \VUgras{spalline}, oppure con un \VUgras{elmo}, una \VUgras{corazza}
\textsuperscript{(3)} e una \VUgras{pistola} \textsuperscript{(4)}. Tutto questo mi ha
interessato molto più delle \VUgras{lanterne magiche} \textsuperscript{(2)}.

%%K t16-tit-1 sub
\VUsaut{3.0mm}
\VUpk{10.2pt}{40}{Belle cose da vedere.}
\VUsaut{2.6mm}

%%K t16-03-1 p
3. --- Nello stesso reparto c'era una \VUgras{sentinella} \textsuperscript{(5)} nella
sua \VUgras{garitta}; sembrava fare la guardia a tutto un serraglio di cartone. Che
quantità di animali! Un \VUgras{pappagallo} \textsuperscript{(6)} sul suo trespolo, un
\VUgras{rinoceronte} \textsuperscript{(7)} con un corno sul naso, una \VUgras{renna}
\textsuperscript{(8)}, uno \VUgras{struzzo} \textsuperscript{(9)}, una
\VUgras{scimmia} \textsuperscript{(10)} che rompe una noce, una \VUgras{volpe}
\textsuperscript{(11)} che porta via una gallina, un \VUgras{coccodrillo}
\textsuperscript{(12)} che spalanca le sue enormi fauci, un \VUgras{cammello}
\textsuperscript{(13)}, un elegante \VUgras{levriero} \textsuperscript{(14)}, una
\VUgras{pantera} \textsuperscript{(15)}, e poi due bestie che non conoscevo e che sono,
mi hanno detto, una \VUgras{donnola} \textsuperscript{(16)} e una \VUgras{iena}
\textsuperscript{(17)}. C'erano anche un \VUgras{leopardo} \textsuperscript{(18)} e un
\VUgras{lupo} \textsuperscript{(21)}; lì accanto c'erano carini \VUgras{topi}
\textsuperscript{(19)} e minuscoli \VUgras{topolini} \textsuperscript{(20)} di gomma.
Infine un \VUgras{ippopotamo} \textsuperscript{(22)} e un \VUgras{dromedario}
\textsuperscript{(23)} con la sua gobba completavano la raccolta. Sarei rimasto lì per
ore, se lì vicino uno splendido campo pieno di \VUgras{soldatini di piombo}
\textsuperscript{(29)} non mi avesse attirato l'occhio.

%%K t16-tit-2 sub
\VUsaut{4.0mm}
\VUpk{10.2pt}{40}{Soldati e guerra.}
\VUsaut{2.8mm}

%%K t16-04-1 p
4. --- Mio nonno, che è un \VUgras{invalido di guerra} \textsuperscript{(24)} e che ha
dovuto lasciare l'esercito dopo una \VUgras{ferita} che gli è valsa la
\VUgras{medaglia al valor militare}, ama raccontarmi le \VUgras{campagne} a cui ha
partecipato. Era felice di spiegarmi i vari movimenti del piccolo esercito in
miniatura. Un gruppo di soldati veri che visitavano il bazar --- un \VUgras{geniere}
\textsuperscript{(25)}, un \VUgras{alpino} \textsuperscript{(26)} con il suo
\VUgras{cappello}, un \VUgras{sergente maggiore} \textsuperscript{(27)} di fanteria e
un \VUgras{sergente} \textsuperscript{(28)} d'artiglieria con un \VUgras{gallone} d'oro
sulla manica --- si è fermato vicino a noi per ascoltare le spiegazioni.

%%K t16-05-1 p
5. --- Mio nonno, molto fiero di avere un pubblico così illustre, ha improvvisato tutto
un piano di battaglia.

%%K t16-05-2 p
La città fortificata, con le sue \VUgras{mura} e i suoi \VUgras{fossati}, era assediata
dall'\VUgras{esercito nemico}, che dopo un lungo \VUgras{bombardamento} si preparava a
darle l'\VUgras{assalto}. Ma gli \VUgras{assediati}, spinti dalla fame, avevano tentato
una \VUgras{sortita} contro il \VUgras{campo} degli \VUgras{assedianti}. Dopo aver
respinto l'\VUgras{attacco} in un \VUgras{combattimento} accanito intorno alle
\VUgras{tende}, gli assedianti \VUgras{vittoriosi} hanno mandato i loro
\VUgras{squadroni} all'\VUgras{inseguimento} degli assalitori \VUgras{sconfitti}.
Questi si sono \VUgras{ritirati}, lasciando il \VUgras{campo di battaglia} coperto di
\VUgras{morti} e di \VUgras{feriti}. I \VUgras{portaferiti} hanno portato i feriti
all'\VUgras{ospedale da campo} perché il \VUgras{chirurgo} li curasse.

%%K t16-05-3 p
Nonostante l'eroica resistenza di alcuni \VUgras{battaglioni} che si erano messi in
\VUgras{quadrato}, la \VUgras{sconfitta} è stata completa; il resto dell'esercito è
\VUgras{fuggito}, lasciando molti \VUgras{prigionieri}. Il nemico stava per coronare la
sua \VUgras{vittoria} prendendo la città. Forse sarebbe stata la fine della campagna, e
dopo un \VUgras{armistizio} si sarebbe potuto firmare un \VUgras{trattato di pace}.

%%K t16-tit-3 sub
\VUsaut{4.4mm}
\VUpk{10.2pt}{40}{I regali di Capodanno.}
\VUsaut{2.8mm}

%%K t16-06-1 p
6. --- I giovani soldati, che ridevano ascoltando il racconto colorito del veterano,
gli hanno fatto ancora qualche domanda. Io intanto ho girato intorno al banco per
guardare una bella \VUgras{balestra} \textsuperscript{(32)} e un \VUgras{arco}
\textsuperscript{(33)} con la sua \VUgras{freccia} \textsuperscript{(31)}. Lì ho
trovato un'altra raccolta di animali: due \VUgras{uccellini canterini}
\textsuperscript{(34)}, un \VUgras{usignolo} e una \VUgras{capinera} a molla che
cantavano a squarciagola, e una serie di bestie strane --- un \VUgras{pipistrello}
\textsuperscript{(35)}, un \VUgras{boa} \textsuperscript{(36)}, una \VUgras{tartaruga}
\textsuperscript{(37)}, una \VUgras{foca} \textsuperscript{(38)}, una
\VUgras{balena} \textsuperscript{(39)} e uno \VUgras{squalo} \textsuperscript{(40)} di
pelle di batti-oro.

%%K t16-07-1 p
7. --- Non mi sono fermato a lungo a guardarli, perché avevo intravisto quello che
sognavo: un magnifico \VUgras{teatrino delle marionette} \textsuperscript{(41)} con
bellissime \VUgras{scene} e un delizioso \VUgras{sipario} rosso. Sul \VUgras{palco} due
attori sembravano recitare una \VUgras{parte} in una \VUgras{commedia} molto
interessante, perché i piccoli \VUgras{spettatori} di cartone nei palchi, o seduti in
\VUgras{platea} e nel \VUgras{loggione} dietro il \VUgras{direttore d'orchestra},
applaudivano con calore.

%%K t16-07-2 p
Purtroppo quel teatro costava molto caro.

%%K t16-08-1 p
8. --- Del resto scegliere fra i giocattoli era difficile, perché le vetrine erano piene
di giocattoli di ogni genere: \VUgras{Arlecchino} \textsuperscript{(42)},
\VUgras{Pulcinella} \textsuperscript{(43)}, \VUgras{Pierrot} \textsuperscript{(44)},
\VUgras{gufi} \textsuperscript{(45)} che battono le ali, \VUgras{merli}
\textsuperscript{(46)} che fischiano, \VUgras{barboncini} che abbaiano, un
\VUgras{grammofono} \textsuperscript{(47)} e \VUgras{rompicapi}. Uno di questi
giocattoli mi è piaciuto moltissimo: rappresentava una \VUgras{battaglia navale}
\textsuperscript{(48)} in cui una \VUgras{nave} è attaccata dai \VUgras{pirati} davanti
a una costa piantata di \VUgras{palme da cocco}.

%%K t16-noto-1 noto
\VUnotes{143.2mm}{%
\textit{(*) Vedi la tavola n. 6.}%
}

%%K t16-09-1 p
9. --- Ho potuto guardare tutte queste meraviglie con tutta calma, perché nel negozio
c'era poca gente --- appena una manciata di sfaccendati, un \VUgras{dragone}
\textsuperscript{(49)} e un \VUgras{esattore} \textsuperscript{(50)} che presentava una
\VUgras{cambiale} alla \VUgras{cassa} centrale \textsuperscript{(51)}.

%%K t16-10-1 p
10. --- Ho chiesto a mio nonno a che cosa servissero i libri sugli scaffali dietro la
\VUgras{cassiera}; mi ha detto che erano i \VUgras{libri contabili}.

%%K t16-tit-4 sub
\VUsaut{3.6mm}
\VUpk{10.2pt}{40}{Uno sguardo per il negozio.}
\VUsaut{2.8mm}

%%K t16-11-1 p
11. --- Lasciato il reparto dei giocattoli, siamo andati alla selleria a dare
un'occhiata alle \VUgras{selle} \textsuperscript{(53)}, alle \VUgras{staffe}
\textsuperscript{(54)}, alle \VUgras{briglie} \textsuperscript{(55)}, ai
\VUgras{morsi} \textsuperscript{(56)} e ai \VUgras{frustini}. Mentre il
\VUgras{capo reparto} \textsuperscript{(57)} dava qualche informazione a mio nonno, io
sono andato a guardare l'esposizione di \VUgras{porcellane} e di
\VUgras{cristallerie} \textsuperscript{(58)}, dove c'erano bellissime
\VUgras{insalatiere} e delicate \VUgras{teiere} \textsuperscript{(59)}.

%%K t16-12-1 p
12. --- Per salire al primo piano siamo passati dietro la vetrina dei \VUgras{busti}
\textsuperscript{(60)}, accanto al reparto della maglieria, dove erano esposte
finissime \VUgras{mutande} \textsuperscript{(63)}. Lì vicino un \VUgras{commesso}
\textsuperscript{(61)} aiutava una \VUgras{cliente} \textsuperscript{(62)} a scegliere
belle \VUgras{stoffe} di seta \textsuperscript{(64)}.

%%K t16-12-2 p
Salendo le scale mi sono accorto che il negozio era decorato con \VUgras{lanterne}
\textsuperscript{(65)} giapponesi, \VUgras{parasoli} \textsuperscript{(66)} ed enormi
\VUgras{palme} \textsuperscript{(70)}.

%%K t16-13-1 p
13. --- Al primo piano non c'era molto da vedere: solo mobili (*), \VUgras{casseforti}
\textsuperscript{(67)}, \VUgras{cassettoni} \textsuperscript{(68)} e
\VUgras{carrozzine} \textsuperscript{(69)}. Ma la vista d'insieme del negozio era molto
\VUgras{divertente}.

%%K t16-14-1 p
14. --- Sotto di noi vedevo gli strumenti musicali: \VUgras{arpe}
\textsuperscript{(71)}, \VUgras{chitarre} \textsuperscript{(72)},
\VUgras{mandolini} \textsuperscript{(73)}, \VUgras{cornette} \textsuperscript{(74)},
\VUgras{clarinetti} \textsuperscript{(75)}, violini con i loro \VUgras{archi}
\textsuperscript{(76)} e perfino una \VUgras{grancassa} \textsuperscript{(77)}. Poco
più in là, al banco della cartoleria, uno \VUgras{studente} \textsuperscript{(78)}
contrattava con il \VUgras{commesso} \textsuperscript{(79)} il prezzo di una cartella.
Vicino a loro riuscivo a distinguere un bell'\VUgras{abbecedario}
\textsuperscript{(80)}.

%%K t16-14-2 p
In mezzo al negozio c'era un alto \VUgras{albero di Natale} \textsuperscript{(81)}
tutto carico di candele, di ninnoli e di giocattoli di ogni genere:
\VUgras{cavallini a dondolo} \textsuperscript{(82)}, \VUgras{bambole}
\textsuperscript{(83)} e così via.

%%K t16-14-3 p
La \VUgras{profumeria} \textsuperscript{(84)} con i suoi \VUgras{collutori} e i suoi
vari \VUgras{profumi} mandava un odore molto gradevole che saliva fino a noi.

%%K t16-tit-5 sub
\VUsaut{3.4mm}
\VUpk{10.2pt}{40}{Compriamo qualcosa.}
\VUsaut{2.8mm}

%%K t16-15-1 p
15. --- Siccome mio nonno cominciava ad annoiarsi, siamo scesi per il grande scalone.
Si è fermato un momento davanti a una \VUgras{carta astronomica}
\textsuperscript{(85)} che mostrava, mi ha detto, le \VUgras{comete}
\textsuperscript{(a)}, le \VUgras{eclissi di luna e di sole} \textsuperscript{(b)}, una
rosa dei venti con i \VUgras{quattro punti cardinali} \textsuperscript{(c)}
\VUgras{(nord, sud, est e ovest)}, una carta dei due \VUgras{emisferi}
\textsuperscript{(d)}, i \VUgras{pianeti} \textsuperscript{(e)} e il
\VUgras{sistema solare} \textsuperscript{(f)}. Io non ci ho capito niente, e mi
interessavano molto di più la livrea gallonata di un \VUgras{fattorino}
\textsuperscript{(86)} che passava e i graziosi \VUgras{ventagli}
\textsuperscript{(87)} accanto a me, vicino ai \VUgras{portamonete}
\textsuperscript{(88)} e ai \VUgras{portafogli} \textsuperscript{(89)}.

%%K t16-16-1 p
16. --- Sul banco ho ammirato anche certe \VUgras{piume} \textsuperscript{(90)} e certe
\VUgras{fibbie da cintura} \textsuperscript{(91)}, e i graziosi \VUgras{fiori finti}
\textsuperscript{(94)} disposti con eleganza nei cesti. Le \VUgras{violette}, le
\VUgras{violacciocche}, i \VUgras{giacinti} \textsuperscript{(95)}, i
\VUgras{gerani} \textsuperscript{(96)}, i \VUgras{gigli} \textsuperscript{(97)} e i
\VUgras{nasturzi} \textsuperscript{(98)} erano imitati benissimo.

%%K t16-tit-6 sub
\VUsaut{4.4mm}
\VUpk{10.2pt}{40}{Il regalo del nonno.}
\VUsaut{2.8mm}

%%K t16-17-1 p
17. --- Ho chiesto a mio nonno di comprarmi uno degli \VUgras{spazzolini da denti}
\textsuperscript{(92)} di una scatola accanto alle \VUgras{forcine}
\textsuperscript{(93)}. Ha accettato, e sono andato io stesso a pagare il mio acquisto
alla cassa, accanto alla quale si stava preparando un grosso \VUgras{pacco}
\textsuperscript{(100)} per una \VUgras{cliente} \textsuperscript{(99)}.

%%K t16-18-1 p
18. --- All'improvviso c'è stato un piccolo trambusto fra le persone che stavano vicino
ai \VUgras{bronzi} \textsuperscript{(101)} artistici. Un \VUgras{ladro}
\textsuperscript{(102)} era appena stato colto sul fatto da un
\VUgras{sorvegliante del negozio} \textsuperscript{(103)} mentre cercava di derubare
qualcuno. Hanno chiamato una guardia per \VUgras{arrestarlo}, e proprio mentre
uscivamo il colpevole veniva portato via al posto di polizia.
\VUsaut{6.0mm}
\VUfilet{22mm}
